\section{Metaheur\'istica}

% TODO (20 pts implementación): cuál elegiste (recocido simulado / tabú /
% genético) y por qué. Describe el estado, la función de vecindad/
% perturbación, y el criterio de aceptación (o el esquema de enfriamiento,
% si es recocido simulado).

\subsection{N-reinas}
% TODO (5 pts aplicar, junto con backtracking): función objetivo (p.ej.
% número de pares de reinas atacándose) y resultados N=8, N=100. Compara
% tiempo y calidad contra backtracking (results/tables/nqueens.csv).

\subsection{Coloreado de grafos}
% TODO (5 pts aplicar): función objetivo (p.ej. número de aristas en
% conflicto) y resultados en 50 y 1000 nodos. Compara contra backtracking
% (results/tables/graph_coloring.csv).

Función objetivo: número de aristas cuyos dos extremos comparten color (\code{count\_conflicts}
en \code{src/problems/graph\_coloring.py}) --- 0 significa coloración válida con exactamente $k$
colores. Igual que en backtracking, se busca incrementalmente el menor $k$ donde el recocido
simulado converge a 0 conflictos.

\textbf{50 nodos}: converge a $k=15$ en 0.24s (un color más que el óptimo de backtracking/
OR-Tools, $k=14$) --- consistente con ser una heurística, no un método exacto: encuentra el
primer $k$ factible de su búsqueda, no necesariamente el mínimo global.

\textbf{1000 nodos, $k=302$}: se probaron los tres esquemas de enfriamiento (geométrico,
exponencial, sinusoidal) con 1800s cada uno. \textbf{Ninguno llega a 0 conflictos} ---
geométrico y exponencial terminan en 259 conflictos ($\sim$60{,}000 iteraciones cada uno),
sinusoidal en 233 (el mejor de los tres, $\sim$60{,}200 iteraciones). Los parámetros por defecto
(\code{initial\_temperature=10}, \code{cooling\_rate=0.995}) están calibrados para instancias
mucho más chicas/dispersas que las probadas en las Secciones~2--3; a esta densidad (0.90) y
dominio (302 colores), la metaheurística tampoco es competitiva --- a diferencia de backtracking
(que ni siquiera avanza, Secci\'on~\ref{sec:coloring-backtracking-1000}), aquí sí hay progreso
medible pero insuficiente en el presupuesto disponible.

\subsection{Comparaci\'on con OR-Tools}
\label{sec:coloring-ortools-1000}
% TODO (5 pts): tabla de results/tables/coloring_comparison.csv --- tiempo
% de cómputo y número de colores usado por backtracking, tu metaheurística,
% y OR-Tools (CP-SAT) en las mismas instancias. Discute las diferencias.

\begin{table}[H]
\centering
\begin{tabular}{@{}lccc@{}}
\toprule
Instancia & Backtracking & Recocido simulado & OR-Tools (CP-SAT) \\
\midrule
50 nodos & $k=14$ (2.41s) & $k=15$ (0.24s) & $k=14$ (0.20s) \\
1000 nodos & no resuelve (1800s, 2 nodos) & no converge (1800s, 233 confl.) & $k=302$ (858s, \textbf{óptimo probado}) \\
\bottomrule
\end{tabular}
\caption{Coloreado de grafos: colores usados y tiempo por método. Para 1000 nodos, $k=302$ es
el mejor valor confirmado factible mediante búsqueda binaria/galopante sobre OR-Tools (ver
\code{experiments/graph\_coloring\_gallop\_overnight.py}); el número cromático exacto no se
determinó por completo dentro del presupuesto de cómputo disponible (ventana final $[48, 302]$).}
\end{table}

En 50 nodos los tres métodos coinciden o casi coinciden con el óptimo real ($k=14$), y OR-Tools
es el más rápido incluso habiendo probado $k=10$ a $13$ (todos infactibles o indecisos) antes de
confirmar $k=14$. En 1000 nodos la diferencia es cualitativa, no solo de velocidad: OR-Tools es
el \'unico de los tres capaz de resolver algo --- backtracking ni siquiera avanza por el costo
por nodo del dominio ancho, y la metaheurística progresa pero no converge. Esto refleja que
CP-SAT usa propagación de restricciones y aprendizaje de cláusulas mucho más sofisticados que
nuestra combinación de FC+AC3+MRV+LCV, una ventaja que se vuelve decisiva conforme crece la
densidad y el tamaño de la instancia. El número cromático exacto de \code{gc\_1000\_9.txt}
sigue sin determinarse: se acotó a $[48, 302]$ (clique de tamaño 49 como cota inferior probada,
$k=302$ como cota superior probada) mediante búsqueda binaria/galopante con OR-Tools, mucho más
efectiva para este propósito que backtracking o la metaheurística por la misma razón anterior
--- ver \code{experiments/graph\_coloring\_gallop\_overnight.py} y
\code{results/tables/graph\_coloring\_gallop\_overnight.csv} para el detalle de la búsqueda.
